\section{Related Work}

\subsection{Long-horizon language agents.}
Language-model agents are commonly instantiated as interleaved reasoning-and-acting systems that issue tool calls, observe environment feedback, and revise their plans over multiple steps \citep{yao2023react}. Recent benchmarks have pushed this setting toward realistic long-horizon tasks, including command-line environments \citep{merrill2026terminalbench}, tool-agent-user interactions with domain rules \citep{yao2024taubench}, machine-learning engineering \citep{chan2024mlebench}, etc. These settings expose failures that are not captured by single-turn reasoning benchmarks: agents may repeat failed actions, lose track of requirements, or drift away from earlier discoveries. Our work focuses on this form of behavioral state decay and studies memory as a mechanism for maintaining decision-relevant execution state during such trajectories.

\subsection{Memory for language-model agents.}
A large body of work augments language models with external or long-term memory. Retrieval-augmented models use non-parametric stores to improve factual grounding and knowledge access \citep{guu2020realm,lewis2020rag}, while long-context memory architectures cache or retrieve past representations to extend the effective context available to the model \citep{wang2023longmem}. Agent memory systems further organize past experience into natural-language records for long-term interaction, personalization, and behavior generation \citep{park2023generative,zhong2024memorybank,packer2023memgpt,zhang2024memorysurvey}. Lifelong agents such as Voyager store reusable skills and retrieve them for new tasks \citep{wang2024voyager}, and production memory layers such as Mem0 emphasize persistent user/session/agent memory with efficient storage and retrieval \citep{mem0ai2026mem0}. Our work is complementary to these systems. Rather than optimizing persistent memory storage or retrieval alone, we ask when remembered execution state should be reactivated as an intervention in an ongoing loop.

\subsection{Learned memory policies and context management.}
Long context windows do not by themselves guarantee reliable use of prior information: models can under-utilize information depending on its position in the context \citep{liu2024lost}. This has motivated work on context compression, memory editing, and learnable context curation. Memory-as-Action formulates working-memory management as explicit editing actions optimized jointly with task performance \citep{zhang2025memact}. Context-Folding lets an agent branch into sub-trajectories for subtasks and then fold completed branches into concise summaries, with reinforcement learning rewards for decomposition and context management \citep{sun2025contextfolding}. Mem-$\alpha$ trains agents to manage memory through tool-based operations over sequential information chunks, using downstream question-answering accuracy over the full interaction history as the reward signal \citep{wang2025memalpha}. Our work shares with these methods the view that memory and context management should be optimized for downstream outcomes, but targets a different control problem. We do not train the action agent to curate its own context, optimize memory construction for later question answering, or fold completed sub-tasks into summaries. Instead, a separate memory agent repeatedly observes a live multi-turn task trajectory, maintains structured execution state, and decides whether any part of that state should intervene in the next action-agent call. This creates a more heterogeneous credit-assignment problem: across tasks, memory may need to remain silent, prevent repeated failures, or reactivate diagnoses and open sub-goals, and unnecessary interventions can be harmful rather than merely redundant.

\subsection{Reflection, critics, and advisor models.}
Another line of work improves agents through feedback, reflection, or auxiliary reasoning. Self-Refine iteratively improves model outputs using self-generated feedback \citep{madaan2023selfrefine}; Reflexion stores verbal reflections from prior attempts in an episodic memory buffer \citep{shinn2023reflexion}; and Tree of Thoughts expands inference-time reasoning through search over intermediate thoughts \citep{yao2023tree}. More recently, advisor-style methods use a second model to steer a black-box or executor model with natural-language guidance \citep{asawa2025advisor,anthropic2026advisor}. Our memory agent is related to these supervisory approaches, but is more constrained: it is not intended to provide broad strategic advice or perform general deliberation. Its output is a memory-grounded reminder derived from maintained execution state, such as a task requirement, environment fact, failed attempt, diagnosis, or open subgoal that is at risk of becoming behaviorally inactive.

\subsection{Positioning.}
Taken together, prior work shows that memory, retrieval, reflection, context compression, and auxiliary guidance can improve language agents. Our contribution is to formulate memory as a selective intervention policy for long-horizon execution. The central question is not only what should be remembered or summarized, but whether, when, and how remembered execution state should enter the action agent's context so that it changes the next decision without unnecessary token or latency overhead. Because tasks differ substantially in environment dynamics, observability, and failure modes, this requires intervention calibration rather than a fixed summarization policy.
